% Source/coverage: Weinberg QFT Vol. I, Ch. 8, Sec. 8.8, physical PDF pp. 392--395
% (printed pp. 369--372, ending immediately before Appendix Traces); equations (8.8.1)--(8.8.15).
% Four ordinary notes including one title note; source-reviewed against page PNGs.

\subsection[Generalization: \texorpdfstring{$p$}{p}-Form Gauge Fields]{Generalization: \texorpdfstring{$p$}{p}-Form Gauge
Fields\texorpdfstring{\protect\footnotemark}{}}
\footnotetext{This section lies somewhat out of the book's main line of
development, and may be omitted in a first reading.}

The antisymmetric field strength tensor $F_{\mu\nu}$ of electromagnetism is a
special case of a general class of tensors of special importance in physics
and mathematics. A $p$-form is an antisymmetric covariant tensor of rank $p$.
From a $p$-form $t_{\mu_1\mu_2\cdots\mu_p}$ one may construct a $(p+1)$-form
called the \emph{exterior derivative}\footnote{Exterior derivatives and
$p$-forms play a special role in general relativity, in part because the
exterior derivative of a tensor transforms like a tensor even though it is
calculated using ordinary rather than covariant derivatives~\hyperref[ref:8.5]{[5]}.} $\dd t$ by
taking the derivative and then antisymmetrizing with respect to all indices:
\begin{align}
(\dd t)_{\mu_1\mu_2\cdots\mu_{p+1}}
&=\partial_{[\mu_1}t_{\mu_2\mu_3\cdots\mu_{p+1}]}
\notag\\
&\equiv\partial_{\mu_1}t_{\mu_2\mu_3\cdots\mu_{p+1}}
-\partial_{\mu_2}t_{\mu_1\mu_3\cdots\mu_{p+1}}+\cdots
+(-1)^p\partial_{\mu_{p+1}}t_{\mu_1\mu_2\cdots\mu_p},
\tag{8.8.1}\label{eq:8.8.1}
\end{align}
with square brackets indicating antisymmetrization with respect to the
indices within the brackets. Because derivatives commute, repeated exterior
derivatives vanish
\begin{equation}
d(\dd t)=0.
\tag{8.8.2}\label{eq:8.8.2}
\end{equation}
A $p$-form whose exterior derivative vanishes is called \emph{closed}, while
a $p$-form that is itself an exterior derivative is called \emph{exact}. From
Eq.~\eqref{eq:8.8.2} it follows that any exact $p$-form is closed; a famous theorem of
Poincaré~\hyperref[ref:8.6]{[6]} states that in a simply connected region, any closed $p$-form is
exact.\footnote{In multiply connected spaces closed forms are not necessarily
exact; although it is possible to write a closed $p$-form as an exterior
derivative locally, this cannot in general be done smoothly throughout the
space. The set of closed $p$-forms, modulo exact $p$-forms, makes up what is
called the $p$th de Rham cohomology group of the space. There is a deep
relation between the de Rham cohomology groups of a space and its topology,
which will be discussed further in Volume II.} For instance, the homogeneous
Maxwell equations~\eqref{eq:8.1.16} tell us that the electromagnetic field strength
two-form $F_{\mu\nu}$ is closed; Poincaré's theorem then shows that it is also
exact, so that it can be written as an exterior derivative, i.e., as
$F_{\mu\nu}=\partial_\mu A_\nu-\partial_\nu A_\mu$. Again using Eq.~\eqref{eq:8.8.2},
we see that the two-form $F_{\mu\nu}$ is invariant if $A_\mu$ is changed by
an exterior derivative, i.e., by a gauge transformation
$\delta A_\mu=\partial_\mu\Omega$.

The formalism of $p$-forms and exterior derivatives makes it natural to
consider the possibility of massless particles described by $p$-form field
gauge fields\footnote{We are speaking loosely in calling
$A_{\mu_1\cdots\mu_p}$ a $p$-form, because for $F=\dd A$ to be a tensor it is
only necessary for $A$ to be a tensor up to a gauge transformation. In fact,
we have already seen that in four spacetime dimensions, it is not possible to
construct a four-vector field from the creation and annihilation operators of
physical massless particles of helicity $\pm1$, so we have to deal with an
$A^\mu(x)$ that according to Eq.~\eqref{eq:8.1.2} transforms as a four-vector only up
to a gauge transformation.} $A_{\mu_1\cdots\mu_p}$, with an invariance under
gauge transformations
\begin{equation}
\delta A=\dd\Omega
\tag{8.8.3}\label{eq:8.8.3}
\end{equation}
or in more detail
\[
\delta A_{\mu_1\cdots\mu_p}
=\partial_{[\mu_1}\Omega_{\mu_2\cdots\mu_p]},
\]
where $\Omega_{\mu_1\cdots\mu_{p-1}}$ is an arbitrary $(p-1)$-form. From such
a $p$-form gauge field, we can construct a gauge-invariant field strength
tensor given by
\begin{equation}
F=\dd A
\tag{8.8.4}\label{eq:8.8.4}
\end{equation}
or in detail
\begin{equation}
F_{\mu_1\cdots\mu_{p+1}}
=\partial_{[\mu_1}A_{\mu_2\cdots\mu_{p+1}]}.
\tag{8.8.5}\label{eq:8.8.5}
\end{equation}
(Alternatively, we can start with a $(p+1)$-form $F$, and from an assumed
condition $\dd F=0$ infer the existence of a $p$-form $A$ with $F=\dd A$.) By
analogy with electrodynamics, we might expect the Lagrangian density for $A$
to take the form
\begin{equation}
\mathcal L=-{1\over2(p+1)}F_{\mu_1\cdots\mu_{p+1}}
F^{\mu_1\cdots\mu_{p+1}}+J^{\mu_1\cdots\mu_p}A_{\mu_1\cdots\mu_p},
\tag{8.8.6}\label{eq:8.8.6}
\end{equation}
where $J$ is an antisymmetric tensor current (either a $c$-number, or a
function of fields other than $A$) that in order to make the action
gauge-invariant must satisfy the conservation condition
\begin{equation}
\partial_{\mu_1}J^{\mu_1\cdots\mu_p}=0.
\tag{8.8.7}\label{eq:8.8.7}
\end{equation}
The Euler--Lagrange equations are then
\begin{equation}
\partial_{\mu_1}F^{\mu_1\cdots\mu_{p+1}}=-J^{\mu_2\cdots\mu_{p+1}}.
\tag{8.8.8}\label{eq:8.8.8}
\end{equation}

These $p$-form gauge fields play an important role in theories with more than
four spacetime dimensions. For instance, in the simplest string theories in
26 spacetime dimensions, there is a normal mode of the string represented at
low energies by a two-form gauge field $A_{\mu\nu}$. But in four spacetime
dimensions, $p$-forms offer no new possibilities.

To see this, note first that in $D$ spacetime dimensions there are no
antisymmetric tensors with more than $D$ indices, so in general we must take
$p+1\leq D$. Like any other $(p+1)$-form with $p+1\leq D$, the field
strength $F$ may be expressed in terms of a dual $(D-p-1)$-form
$\mathcal F$, as
\begin{equation}
F^{\mu_1\cdots\mu_{p+1}}
=\epsilon^{\mu_1\cdots\mu_D}
\mathcal F_{\mu_{p+2}\cdots\mu_D}.
\tag{8.8.9}\label{eq:8.8.9}
\end{equation}
Likewise, the $p$-form current $J$ may be expressed in terms of a dual
$(D-p)$-form current $\mathcal J$, as
\begin{equation}
J^{\mu_1\cdots\mu_p}
=\epsilon^{\mu_1\cdots\mu_D}\mathcal J_{\mu_{p+1}\cdots\mu_D}.
\tag{8.8.10}\label{eq:8.8.10}
\end{equation}
The field equation~\eqref{eq:8.8.8} and conservation condition \eqref{eq:8.8.7} then read
simply
\begin{equation}
\dd\mathcal F=\mathcal J,\qquad \dd\mathcal J=0.
\tag{8.8.11}\label{eq:8.8.11}
\end{equation}
Because the dual current $\mathcal J$ is closed, it may be written in terms
of a $(D-p-1)$-form $\mathcal S$ as
\begin{equation}
\mathcal J=\dd\mathcal S.
\tag{8.8.12}\label{eq:8.8.12}
\end{equation}
Eqs.~\eqref{eq:8.8.11} and \eqref{eq:8.8.12} tell us that the difference of $\mathcal F$ and
$\mathcal S$ is closed, and therefore according to Poincaré's theorem, may be
written as
\begin{equation}
\mathcal F=\mathcal S+\dd\phi,
\tag{8.8.13}\label{eq:8.8.13}
\end{equation}
with $\phi$ a $(D-p-2)$-form. There is an exception for the case $p=D-1$,
where $\mathcal F$ and $\mathcal S$ are zero-forms, i.e., scalars, and the
condition $\dd\mathcal F=\dd\mathcal S$ simply tells us that $\mathcal F$ and
$\mathcal S$ differ only by a constant. In this case the gauge field describes
no degrees of freedom at all. We may therefore concern ourselves only with
the cases $p\leq D-2$.

For $p\leq D-2$, the homogeneous 'Maxwell' equations $\dd F=0$ read
\begin{equation}
\partial_{\mu_1}\mathcal F^{\mu_1\cdots\mu_{D-p}}=0,
\tag{8.8.14}\label{eq:8.8.14}
\end{equation}
which with Eq.~\eqref{eq:8.8.13} yields the field equation for $\phi$:
\begin{equation}
\partial_{\mu_1}(\dd\phi)^{\mu_1\cdots\mu_{D-p}}
=-(\dd\mathcal S)^{\mu_2\cdots\mu_{D-p}}.
\tag{8.8.15}\label{eq:8.8.15}
\end{equation}
This is invariant under a new set of gauge transformations
$\phi\to\phi+\dd\omega$, except that where $D-p-2=0$ the gauge transformation
that leaves $F$ invariant is $\phi\to\phi+c$, where $c$ is an arbitrary
constant. \emph{We see that in $D$ spacetime dimensions, the theory of a
$p$-form gauge field $A$ with current $J$ is equivalent to the theory of a
$(D-p-2)$-form gauge field $\phi$, with current $-\dd\mathcal S$.}

We can now understand why $p$-form gauge fields offer no new possibilities in
four spacetime dimensions. As we have seen, we need only consider the cases
$p\leq D-2$, or $p=0$, 1, or 2. A zero-form gauge field is a scalar $S$, for
which Eq.~\eqref{eq:8.8.5} reads $F_\mu=\partial_\mu S$, and the field equations
\eqref{eq:8.8.8} read simply $\Box S=-J$. The gauge invariance here is invariance
under a shift $S\to S+c$, with $c$ a constant. This is just the theory of a
massless scalar field with only derivative interactions. A one-form gauge
field is a four-vector $A^\mu(x)$ coupled to a conserved four-vector current,
just as in electrodynamics. Finally, according to the general result quoted
above, a two-form gauge field in four spacetime dimensions is equivalent to
a zero-form gauge field, which as we have seen is equivalent to a massless
derivatively coupled scalar field.
